\BeleskeID{06-s022}
% element: 06-s022:e01
% element: 06-s022:d01
Za parnu parnost zaštitni bit jednak je EXOR-u bita poruke. Ako je broj jedinica u poruci neparan, dodaje se jedinica; ako je paran, dodaje se nula. Na prijemu EXOR svih bita, uključujući zaštitni, treba da bude nula.

Svaka promena jednog bita obrće ukupnu parnost. Zato neparan broj promena daje signal greške, dok paran broj promena može proći neprimećeno. Na primer, u 0000 promene dva donja bita daju 0011, koji i dalje ima parnu parnost. Sam zaštitni bit ne određuje poziciju greške i ne omogućava njeno ispravljanje.

Fizički procesi mogu izazvati grešku pri obradi, smeštanju ili prenosu. Dodavanje redundantnog sadržaja omogućava detekciju, a odgovarajući složeniji kodovi i korekciju.
